% ======================================================
% 第 51 天
% ======================================================
\setday{51}{曲面及其方程}

\oneproblem{
    以下方程描述的曲面不是直纹面的是（\quad）
    \begin{multicols}{2}
    \begin{enumerate}[label=(\Alph*)]
        \item $z=xy$
        \item $\frac{x^2}{2}+\frac{y^2}{3}-\frac{z^2}{4}=1$
        \item $x^2+y^2=z^2$
        \item $\frac{x^2}{2}+\frac{y^2}{3}+\frac{z^2}{4}=1$
    \end{enumerate}
    \end{multicols}
}

\oneproblem{
    求直线 $L:\frac{x-1}{1}=\frac{y}{1}=\frac{z-1}{-1}$ 在平面 $\pi:x-y+2z-1=0$ 上的投影直线 $L_0$ 的方程，并求 $L_0$ 绕 $y$ 轴旋转一周所成的曲面方程.
}

\oneproblem{
    椭球面 $S_1$ 是椭圆 $\frac{x^2}{4}+\frac{y^2}{3}=1$ 绕 $x$ 轴旋转而成，圆锥面 $S_2$ 是由过点 $(4,0)$ 且与椭圆 $\frac{x^2}{4}+\frac{y^2}{3}=1$ 相切的直线绕 $x$ 轴旋转而成.
    \begin{enumerate}[label=(\arabic*)]
        \item 求 $S_1$ 及 $S_2$ 的方程；
        \item 求 $S_1$ 与 $S_2$ 之间的立体体积.
    \end{enumerate}
}

\oneproblem{
    设 $N(0,0,1)$ 是球面 $S:x^2+y^2+z^2=1$ 的北极点. $A(a_1,a_2,0),B(b_1,b_2,0),C(c_1,c_2,0)$ 为 $xOy$ 面上不同的三点. 设连接 $N$ 与 $A,B,C$ 的三直线依次交球面 $S$ 于点 $A_1,B_1,C_1$.
    \begin{enumerate}[label=(\arabic*)]
        \item 求连接 $N$ 与 $A$ 两点的直线方程；
        \item 求点 $A_1,B_1,C_1$ 三点的坐标；
        \item 给定点 $A(1,-1,0),B(-1,1,0),C(1,1,0)$，求四面体 $NA_1B_1C_1$ 的体积.
    \end{enumerate}
}

\oneproblem{
    一动直线 $L$ 沿三条直线 $L_1:\frac{x}{2}=\frac{y-1}{0}=\frac{z}{-1}$, $L_2:\frac{x-2}{0}=\frac{y}{2}=\frac{z}{1}$, $L_3:\frac{x}{2}=\frac{y+1}{0}=\frac{z}{1}$ 滑动，求动直线 $L$ 的轨迹方程.
}

\oneproblem{
    有两条互相垂直相交的直线 $L_1$ 和 $L_2$，其中 $L_1$ 绕 $L_2$ 作螺旋运动，即 $L_1$ 一方面绕 $L_2$ 作等速转动，另一方面又沿着 $L_2$ 作等速直线运动，在运动中 $L_1$ 保持与 $L_2$ 垂直相交，这样由 $L_1$ 滑过所形成的轨迹曲面称为螺旋面. 试建立螺旋面的方程.
}

\oneproblem{
    设 $M$ 是以三个正半轴为母线的半圆锥面，求其方程.
}

\oneproblem{
    设 $\Gamma$ 为椭圆抛物面 $z=3x^2+4y^2+1$，从原点作 $\Gamma$ 的切锥面. 求切锥面的方程.
}

\oneproblem{
    已知四点 $(1,2,7),(4,3,3),(5,-1,6),(\sqrt{7},\sqrt{7},0)$. 试求过这四点的球面方程.
}

\oneproblem{
    已知二次曲面 $\Sigma$（非退化）过以下九点：$A(1,0,0),B(1,1,2),C(1,-1,-2),D(3,0,0),E(3,1,2),F(3,-2,-4)$\\$,G(0,1,4),H(3,-1,-2),I(5,2\sqrt{2},8)$. 问 $\Sigma$ 是哪一类曲面？
}

\oneproblem{
    求经过三平行直线 $L_1:x=y=z$, $L_2:x-1=y=z+1$, $L_3:x=y+1=z-1$ 的圆柱面的方程.
}

\oneproblem{
    考虑单叶双曲面 $S:x^2-y^2+z^2=1$.
    \begin{enumerate}[label=(\arabic*)]
        \item 证明：$S$ 上同一族直母线中任意两条不同的直母线是异面直线；
        \item 设 $S$ 上同一族直母线中的两条直母线分别经过 $M_1(1,1,1)$ 与 $M_2(2,2,1)$ 两点，求这两条直母线的公垂线方程以及这两条直母线之间的距离.
    \end{enumerate}
}

\oneproblem{
    由曲面 $\frac{x^2}{a^2}+\frac{y^2}{b^2}+\frac{z^2}{c^2}=1$ 中心引三条互相垂直的射线，分别交曲面于 $P_1,P_2,P_3$. 设 $|\overrightarrow{OP_1}|=r_1,|\overrightarrow{OP_2}|=r_2,|\overrightarrow{OP_3}|=r_3$. 证明：$\frac{1}{r_1^2}+\frac{1}{r_2^2}+\frac{1}{r_3^2}=\frac{1}{a^2}+\frac{1}{b^2}+\frac{1}{c^2}$.
}

\oneproblem{
    设圆锥面的顶点在原点，底面圆是平面 $x+y+z=3$ 上以点 $(1,1,1)$ 为圆心且以 $1$ 为半径的圆，试求圆锥面的方程.
}

\oneproblem{
    求过单叶双曲面 $\frac{x^2}{4}+\frac{y^2}{2}-2z^2=1$ 和球面 $x^2+y^2+z^2=4$ 的交线且与直线 $\begin{cases}x=0,\\3y+z=0\end{cases}$ 垂直的平面方程.
}

\oneproblem{
    给定 $yOz$ 平面上的圆 $C:y=\sqrt{3}+\cos\theta,z=1+\sin\theta(\theta\in[0,2\pi])$.
    \begin{enumerate}[label=(\arabic*)]
        \item 求 $C$ 绕 $z$ 轴旋转所得到的环面 $S$ 的隐式方程.
        \item 设 $z_0\ge0$，以 $M(0,0,z_0)$ 为顶点的两个锥面 $S_1$ 和 $S_2$ 的半顶角之差为 $\pi/3$，且均与环面 $S$ 相切（每条母线都与环面相切），求 $z_0$ 和 $S_1,S_2$ 的隐式方程.
    \end{enumerate}
}

\oneproblem{
    设 $S$ 为 $\mathbb{R}^3$ 中的抛物面 $z=\frac{1}{2}(x^2+y^2)$, $P=(a,b,c)$ 为 $S$ 外一固定点，满足 $a^2+b^2>2c$. 过 $P$ 作 $S$ 的所有切线. 证明：这些切线的切点落在同一张平面上.
}

% ======================================================
% 第 52 天
% ======================================================
\setday{52}{空间曲线及其方程}

\oneproblem{
    设平面曲线 $L$ 的方程为 $Ax^2+By^2+Cxy+Dx+Ey+F=0$，且通过五个点\\ $P_1(-1,0),P_2(0,-1),P_3(0,1),P_4(2,-1)$ 和 $P_5(2,1)$，试求 $L$ 上任意两点之间的直线距离最大值.
}

\oneproblem{
    求曲线 $\begin{cases}z=y^2,\\x=0\end{cases}$ 绕 $z$ 轴旋转的曲面与平面 $x+y+z=1$ 的交线在 $xOy$ 平面的投影曲线方程.
}

\oneproblem{
    试求曲线 $\Gamma:x=\cos t,y=\sin t,z=1,0\le t\le2\pi$ 在平面 $\pi:x+y+z+2=0$ 上的投影曲线方程.
}

\oneproblem{
    求直线 $L_1:\frac{x-1}{1}=\frac{y}{1}=\frac{z-1}{-1}$ 在平面 $\pi:x-y+2z-1=0$ 上的投影直线 $L_0$ 的方程，并求 $L_0$ 绕 $y$ 轴旋转一周所成曲面的方程.
}

\oneproblem{
    求与两直线 $L_1:y=0,z=c$ 与 $L_2:x=0,z=-c\ (c\neq0)$ 均相交，且与双曲线 $\Gamma:xy+c^2=0,z=0$ 也相交的动直线 $L$ 所产生的曲面方程.
}

\oneproblem{
    有一束平行于直线 $L:x=y=-z$ 的平行光束照射不透明球面 $S:x^2+y^2+z^2=2z$. 求球面在 $xOy$ 平面上留下的阴影部分的边界曲线方程.
}

\oneproblem{
    求过两球面的交线 $\Gamma:\begin{cases}x^2+y^2+z^2=5\\(x-2)^2+(y-1)^2+z^2=1\end{cases}$ 的正圆柱面方程（母线垂直于准线所在平面的柱面）.
}

% ======================================================
% 第 53 天
% ======================================================
\setday{53}{向量值函数的导数与积分}

\oneproblem{
    证明：$\lim_{t\to a}\mathbf{r}(t)=\mathbf{b}$ 的充分必要条件为：对任意的 $\varepsilon>0$，存在实数 $\delta>0$，当 $0<|t-a|<\delta$ 时，$|\mathbf{r}(t)-\mathbf{b}|<\varepsilon$.
}

\oneproblem{
    在军事作战仿真模拟中，若红方导弹和蓝方飞机的运动轨迹分别由向量值函数 $\mathbf{r}_1(t)=(t^2,7t-12,t^2)$, $\mathbf{r}_2(t)=(4t-3,t^2,5t-6)\ (t\ge0)$ 给出，问：
    \begin{enumerate}[label=(\arabic*)]
        \item 当 $t=1$ 时，导弹距离飞机多远？
        \item 导弹能否击中飞机？
    \end{enumerate}
}

\oneproblem{
    证明一元向量值函数中值定理：设向量值函数 $\vec{r}(t)=(f(t),g(t))$ 满足：
    \begin{enumerate}[label=(\arabic*)]
        \item 在闭区间 $[a,b]$ 上连续；
        \item 在开区间 $(a,b)$ 上可导；
        \item $\vec{r}'(t)\neq\vec{0},\forall t\in(a,b)$；
        \item $\vec{r}(a)\neq\vec{r}(b)$；
    \end{enumerate}
    在 $(a,b)$ 内至少存在一点 $\xi$，使得 $\vec{r}'(\xi)=\lambda[\vec{r}(b)-\vec{r}(a)]$，其中 $\lambda=\frac{|\vec{r}'(\xi)|}{|\vec{r}(b)-\vec{r}(a)|}$.
}

\oneproblem{
    在一高为 $400$ 米半顶角为 $\frac{\pi}{6}$ 的圆锥形山包上，敌方火力点位于点 $A$ 处，我方爆破手位于点 $B$ 处，$A$ 点位于 $yOz$ 面的圆锥母线上距顶点 $P$ 处的距离为 $100\sqrt{2}$ 米，$B$ 点位于 $zOx$ 面的圆锥面的母线上距顶点 $P$ 处 $100(1+\sqrt{3})$ 米. 从 $B$ 到 $A$ 的最短距离是将圆锥面沿一条母线展开后的扇形面上 $B$、$A$ 两点的直线距离，试求从 $B$ 到 $A$ 的最短距离曲线的向量方程.
}

\oneproblem{
    求曲线 $\mathbf{r}(t)=t\mathbf{i}-t^2\mathbf{j}+t^3\mathbf{k}$ 上的点，使该点的切线平行于平面 $x+2y+z=4$.
}

\oneproblem{
    设 $\vec{r}(t)$ 是可导的向量值函数，且 $\vec{r}'(t)\neq\vec{0}$，若 $|\vec{r}(t)|=C$（$C$ 为常数），证明：$\vec{r}(t)$ 与 $\vec{r}'(t)$ 垂直.
}

\oneproblem{
    设质量为 $m$ 的质点的位置向量为 $\vec{r}(t)$，它的角动量为 $\vec{L}(t)=m\vec{r}(t)\times\vec{v}(t)$，转动力矩为 $\vec{M}(t)=m\vec{r}(t)\times\vec{a}(t)$. 试证明：$\vec{L}'(t)=\vec{M}(t)$.
}

\oneproblem{
    证明：二阶函数行列式的求导法则：$\frac{d}{dt}\begin{vmatrix}a_1(t)&a_2(t)\\b_1(t)&b_2(t)\end{vmatrix}=\begin{vmatrix}a_1'(t)&a_2'(t)\\b_1(t)&b_2(t)\end{vmatrix}+\begin{vmatrix}a_1(t)&a_2(t)\\b_1'(t)&b_2'(t)\end{vmatrix}$. 并由此导出三阶函数行列式的求导法则 $\frac{d}{dt}\begin{vmatrix}a_1(t)&a_2(t)&a_3(t)\\b_1(t)&b_2(t)&b_3(t)\\c_1(t)&c_2(t)&c_3(t)\end{vmatrix}$.
}

\oneproblem{
    如图，一条河宽 $100$ 米，一艘划艇在时刻 $t=0$ 离开对岸朝着近岸以 $20$ 米/分的速率行使，方向垂直于河岸，水流在 $(x,y)$ 的速度是 $\mathbf{v}=\left[-\frac{1}{250}(y-50)^2+10\right]\mathbf{i}$（米/分），其中 $0<y<100$.
    \begin{enumerate}[label=(\arabic*)]
        \item 设 $\mathbf{r}(0)=100\mathbf{j}$，求划艇的位置向量 $\mathbf{r}(t)$；
        \item 划艇在下游多远处靠近岸？
    \end{enumerate}
}

% ======================================================
% 第 54 天
% ======================================================
\setday{54}{多元函数的极限与连续性}

\oneproblem{
    判定以下函数当 $(x,y)\to(0,0)$ 时极限的存在性.
    \begin{enumerate}[label=(\arabic*)]
        \item $f(x,y)=\frac{x^3y}{x^6+y^2}$；
        \item $\lim_{\substack{x\to0\\y\to0}}\frac{1-\cos(x^2+y^2)}{(x^2+y^2)x^2y^2}$；
        \item $\lim_{\substack{x\to0\\y\to0}}\frac{\sin(x+y)}{x-y}$；
        \item $f(x,y)=x\frac{\ln(1+xy)}{x+y}$.
    \end{enumerate}
}

\oneproblem{
    判断下列二元函数当 $(x,y)\to(0,0)$ 时二重极限与累次极限的存在性：
    \begin{enumerate}[label=(\arabic*)]
        \item $f(x,y)=\frac{xy}{x^2+y^2}$；
        \item $f(x,y)=\frac{x^2-y^2}{x^2+y^2}$；
        \item $f(x,y)=(x+y)\sin\frac{1}{x}\sin\frac{1}{y}$.
    \end{enumerate}
}

\oneproblem{
    讨论函数 $f(x,y)=\begin{cases}\frac{x^2y}{x^2+y^2},&(x,y)\neq(0,0)\\0,&(x,y)=(0,0)\end{cases}$ 的连续性.
}

\oneproblem{
    求下列极限：
    \begin{enumerate}[label=(\arabic*)]
        \item $\lim_{\substack{x\to0\\y\to0}}\frac{x(y-x)}{\sqrt{x^2+y^2}}$.
        \item $\lim_{\substack{x\to\infty\\y\to a}}\left(1+\frac{1}{xy}\right)^{\frac{x^2}{x+y}}\ (a\neq0)$.
        \item $\lim_{\substack{x\to0\\y\to0}}\frac{x^2y}{x^2+y^2}\sin\frac{1}{x^2+y^2}$.
    \end{enumerate}
}

\oneproblem{
    设 $F(x,y)=\frac{f(y-x)}{2x}$ 及 $F(1,y)=\frac{y^2}{2}-y+5$. 又 $x_0>0,x_n=F(x_{n-1},2x_{n-1}),n=1,2,\cdots$，证明：当 $n\to\infty$ 时，数列 $\{x_n\}$ 的极限存在.
}

% ======================================================
% 第 55 天
% ======================================================
\setday{55}{偏导数及其基本计算法}

\oneproblem{
    已知 $f(x,y)=e^{\sqrt{x^2+y^4}}$，则（\quad）
    \begin{multicols}{2}
    \begin{enumerate}[label=(\Alph*)]
        \item $f_x'(0,0),f_y'(0,0)$ 都存在
        \item $f_x'(0,0)$ 不存在，$f_y'(0,0)$ 存在
        \item $f_x'(0,0)$ 不存在，$f_y'(0,0)$ 不存在
        \item $f_x'(0,0),f_y'(0,0)$ 都不存在
    \end{enumerate}
    \end{multicols}
}

\oneproblem{
    考虑二元函数的下面 4 条性质：
    \ding{172} $f(x,y)$ 在点 $(x_0,y_0)$ 处连续，
    \ding{173} $f(x,y)$ 在点 $(x_0,y_0)$ 处的两个偏导数连续，
    \ding{174} $f(x,y)$ 在点 $(x_0,y_0)$ 处可微，
    \ding{175} $f(x,y)$ 在点 $(x_0,y_0)$ 处两个偏导数存在.
    若用“$P\Rightarrow Q$”表示可由性质 $P$ 推出 $Q$，则有（\quad）
    \begin{multicols}{2}
    \begin{enumerate}[label=(\Alph*)]
        \item \ding{173}$\Rightarrow$\ding{174}$\Rightarrow$\ding{172}
        \item \ding{174}$\Rightarrow$\ding{173}$\Rightarrow$\ding{172}
        \item \ding{174}$\Rightarrow$\ding{175}$\Rightarrow$\ding{172}
        \item \ding{174}$\Rightarrow$\ding{172}$\Rightarrow$\ding{175}
    \end{enumerate}
    \end{multicols}
}

\oneproblem{
    设函数 $f(x,y)$ 为可微函数，且对任意的 $x,y$ 都有 $\frac{\partial f}{\partial x}>0,\frac{\partial f}{\partial y}<0$，则使不等式 $f(x_1,y_1)>f(x_2,y_2)$ 成立的一个充分条件是（\quad）
    \begin{multicols}{2}
    \begin{enumerate}[label=(\Alph*)]
        \item $x_1>x_2,y_1<y_2$
        \item $x_1>x_2,y_1>y_2$
        \item $x_1<x_2,y_1<y_2$
        \item $x_1<x_2,y_1>y_2$
    \end{enumerate}
    \end{multicols}
}

\oneproblem{
    关于函数 $f(x,y)=\begin{cases}xy,xy\neq0\\x,y=0\\y,x=0\end{cases}$ 给出下列结论：
    \begin{enumerate}[label=(\arabic*)]
        \item $\left.\frac{\partial f}{\partial x}\right|_{(0,0)}=1$
        \item $\left.\frac{\partial^2 f}{\partial x\partial y}\right|_{(0,0)}=1$
        \item $\lim_{(x,y)\to(0,0)}f(x,y)=0$
        \item $\lim_{y\to0}\lim_{x\to0}f(x,y)=0$
    \end{enumerate}
    其中正确的个数为（\quad）
    \begin{multicols}{2}
    \begin{enumerate}[label=(\Alph*)]
        \item 4
        \item 3
        \item 2
        \item 1
    \end{enumerate}
    \end{multicols}
}

\oneproblem{
    设 $F(x,y)$ 的偏导数存在，$f(x)$ 可导且 $f'(x)\neq0$，求极限 $\lim_{\varepsilon\to0}\frac{F(x,y+\varepsilon)-F(x,y-\varepsilon)}{f(x+\varepsilon)-f(x-\varepsilon)}$.
}

\oneproblem{
    设 $z=\left(\frac{y}{x}\right)^{\frac{x}{y}}$，求 $\left.\frac{\partial z}{\partial x}\right|_{(1,2)}$.
}

\oneproblem{
    已知 $f(x,y)=x^2\arctan\frac{y}{x}-y^2\arctan\frac{x}{y}$，求 $\frac{\partial^2 f}{\partial x\partial y}$.
}

\oneproblem{
    设 $z=\frac{2x}{x^2-y^2}$，求 $\left.\frac{\partial^n z}{\partial y^n}\right|_{(2,1)}$.
}

\oneproblem{
    已知函数 $f(x,y)$ 满足 $\frac{\partial f}{\partial y}=2(y+1)$，且 $f(y,y)=(y+1)^2-(2-y)\ln y$，求曲线 $f(x,y)=0$ 所围图形绕直线 $y=-1$ 旋转所成的旋转体的体积.
}

\oneproblem{
    证明函数 $f(x,y)=\begin{cases}xy\frac{x^2-y^2}{x^2+y^2},&(x,y)\neq(0,0)\\0,&(x,y)=(0,0)\end{cases}$ 在点 $(0,0)$ 处的两个混合偏导数不相等.
}

\oneproblem{
    设 $f_{xy}(x,y),f_{yx}(x,y)$ 都在点 $(x_0,y_0)$ 处连续，证明：$f_{xy}(x_0,y_0)=f_{yx}(x_0,y_0)$.
}

\oneproblem{
    设可微函数 $f(x,y)$ 满足 $\frac{\partial f}{\partial x}=-f(x,y),f\left(0,\frac{\pi}{2}\right)=1$，且 $\lim_{n\to\infty}\left(\frac{f\left(0,y+\frac{1}{n}\right)}{f(0,y)}\right)^n=e^{\cot y}$，试求二元函数 $f(x,y)$ 的表达式.
}

\oneproblem{
    设函数 $f(x,y)$ 的两个偏导数在点 $(x_0,y_0)$ 的某一邻域内存在且有界，证明：$f(x,y)$ 在点 $(x_0,y_0)$ 处连续.
}

\oneproblem{
    设函数 $f(x,y)$ 对每个固定的 $y$ 是变量 $x$ 的连续函数，且有有界的偏导数 $f_y'(x,y)$. 证明：函数 $f(x,y)$ 是连续函数.
}

\oneproblem{
    已知函数 $u(x,y)$ 满足 $2\frac{\partial^2 u}{\partial x^2}-2\frac{\partial^2 u}{\partial y^2}+3\frac{\partial u}{\partial x}+3\frac{\partial u}{\partial y}=0$，求 $a,b$ 的值，使得在变换 $u(x,y)=v(x,y)e^{ax+by}$ 下，上述等式可以化为 $v(x,y)$ 不含一阶偏导数的等式.
}

\oneproblem{
    设函数 $u=f(x,y)$ 具有二阶连续偏导数，且满足等式 $4\frac{\partial^2 u}{\partial x^2}+12\frac{\partial^2 u}{\partial x\partial y}+5\frac{\partial^2 u}{\partial y^2}=0$，确定 $a,b$ 的值，使等式在变换 $\xi=x+ay,\eta=x+by$ 下化简为 $\frac{\partial^2 u}{\partial\xi\partial\eta}=0$.
}

% ======================================================
% 第 56 天
% ======================================================
\setday{56}{全微分及其应用}

\oneproblem{
    二元函数 $f(x,y)$ 在点 $(0,0)$ 处可微的一个充分条件是（\quad）
    \begin{enumerate}[label=(\Alph*)]
        \item $\lim_{(x,y)\to(0,0)}[f(x,y)-f(0,0)]=0$
        \item $\lim_{x\to0}\frac{f(x,0)-f(0,0)}{x}=0$，且 $\lim_{y\to0}\frac{f(0,y)-f(0,0)}{y}=0$
        \item $\lim_{(x,y)\to(0,0)}\frac{f(x,y)-f(0,0)}{\sqrt{x^2+y^2}}=0$
        \item $\lim_{x\to0}[f_x'(x,0)-f_x'(0,0)]=0$，且 $\lim_{y\to0}[f_y'(0,y)-f_y'(0,0)]=0$
    \end{enumerate}
}

\oneproblem{
    如果 $f(x,y)$ 在 $(0,0)$ 处连续，那么下列命题正确的是（\quad）
    \begin{multicols}{2}
    \begin{enumerate}[label=(\Alph*)]
        \item 若极限 $\lim_{\substack{x\to0\\y\to0}}\frac{f(x,y)}{|x|+|y|}$ 存在，则 $f(x,y)$ 在 $(0,0)$ 处可微
        \item 若极限 $\lim_{\substack{x\to0\\y\to0}}\frac{f(x,y)}{x^2+y^2}$ 存在，则 $f(x,y)$ 在 $(0,0)$ 处可微
        \item 若 $f(x,y)$ 在 $(0,0)$ 处可微，则极限 $\lim_{\substack{x\to0\\y\to0}}\frac{f(x,y)}{|x|+|y|}$ 存在
        \item 若 $f(x,y)$ 在 $(0,0)$ 处可微，则极限 $\lim_{\substack{x\to0\\y\to0}}\frac{f(x,y)}{x^2+y^2}$ 存在
    \end{enumerate}
    \end{multicols}
}

\oneproblem{
    函数 $f(x,y)$ 在 $(0,0)$ 处可微，$f(0,0)=0,\vec{n}=\left.\left(\frac{\partial f}{\partial x},\frac{\partial f}{\partial y},-1\right)\right|_{(0,0)}$，非零向量 $\vec{\alpha}$ 与 $\vec{n}$ 垂直，则（\quad）
    \begin{multicols}{2}
    \begin{enumerate}[label=(\Alph*)]
        \item $\lim_{(x,y)\to(0,0)}\frac{|\vec{n}\cdot(x,y,f(x,y))|}{\sqrt{x^2+y^2}}$ 存在
        \item $\lim_{(x,y)\to(0,0)}\frac{|\vec{n}\times(x,y,f(x,y))|}{\sqrt{x^2+y^2}}$ 存在
        \item $\lim_{(x,y)\to(0,0)}\frac{|\vec{\alpha}\cdot(x,y,f(x,y))|}{\sqrt{x^2+y^2}}$ 存在
        \item $\lim_{(x,y)\to(0,0)}\frac{|\vec{\alpha}\times(x,y,f(x,y))|}{\sqrt{x^2+y^2}}$ 存在
    \end{enumerate}
    \end{multicols}
}

\oneproblem{
    已知 $z=\arctan\frac{x+y}{x-y}$，求 $dz$.
}

\oneproblem{
    研究函数 $f(x,y)=\begin{cases}xy\sin\frac{1}{\sqrt{x^2+y^2}},&(x,y)\neq(0,0)\\0,&(x,y)=(0,0)\end{cases}$ 在 $(0,0)$ 处的连续性、偏导数的存在性、偏导数的连续性与可微性.
}

\oneproblem{
    设连续函数 $z=f(x,y)$ 满足 $\lim_{\substack{x\to0\\y\to1}}\frac{f(x,y)-2x+y-2}{\sqrt{x^2+(y-1)^2}}=0$，求 $dz|_{(0,1)}$.
}

\oneproblem{
    设函数 $f(x,y)$ 具有一阶连续偏导数，且 $df(x,y)=ye^y dx+x(1+y)e^y dy,f(0,0)=0$，求函数 $f(x,y)$ 的表达式.
}

\oneproblem{
    已知函数 $z=z(x,y)$ 的全微分为 $dz=[12x^2-y\cos(xy)]dx+[6y-x\cos(xy)]dy$，求 $z(x,y)$ 的函数表达式.
}

\oneproblem{
    设二元函数 $f(x,y)=|x-y|\varphi(x,y)$，其中 $\varphi(x,y)$ 在点 $(0,0)$ 的一个邻域内连续，试证明函数 $f(x,y)$ 在 $(0,0)$ 点处可微的充分必要条件是 $\varphi(0,0)=0$.
}

\oneproblem{
    若函数 $z=f(x,y)$ 的偏导数 $\frac{\partial z}{\partial x},\frac{\partial z}{\partial y}$ 在点 $(x,y)$ 处连续，证明函数 $f(x,y)$ 在 $(x,y)$ 处可微.
}

\oneproblem{
    设 $f(x,y)$ 在区域 $D$ 内可微，且 $\sqrt{\left(\frac{\partial f}{\partial x}\right)^2+\left(\frac{\partial f}{\partial y}\right)^2}\leq M$，$A(x_1,y_1),B(x_2,y_2)$ 是 $D$ 内两点，线段 $AB$ 包含在 $D$ 内. 证明：$|f(x_1,y_1)-f(x_2,y_2)|\leq M|AB|$，其中 $|AB|$ 表示线段 $AB$ 的长度.
}

% ======================================================
% 第 57 天
% ======================================================
\setday{57}{多元复合函数的求导（一）}

\oneproblem{
    设函数 $f(x,y)$ 可微，且 $f(x+1,e^x)=x(x+1)^2,f(x,x^2)=2x^2\ln x$，则 $df(1,1)=$（\quad）
    \begin{multicols}{2}
    \begin{enumerate}[label=(\Alph*)]
        \item $dx+dy$
        \item $dx-dy$
        \item $dy$
        \item $-dy$
    \end{enumerate}
    \end{multicols}
}

\oneproblem{
    设函数 $z=xyf\left(\frac{y}{x}\right)$，且 $f(u)$ 可导，若 $x\frac{\partial z}{\partial x}+y\frac{\partial z}{\partial y}=y^2(\ln y-\ln x)$，则（\quad）
    \begin{multicols}{2}
    \begin{enumerate}[label=(\Alph*)]
        \item $f(1)=\frac{1}{2},f'(1)=0$
        \item $f(1)=0,f'(1)=\frac{1}{2}$
        \item $f(1)=\frac{1}{2},f'(1)=1$
        \item $f(1)=0,f'(1)=1$
    \end{enumerate}
    \end{multicols}
}

\oneproblem{
    设函数 $f(t)$ 连续，令 $F(x,y)=\int_0^{x-y}(x-y-t)f(t)dt$，则（\quad）
    \begin{multicols}{2}
    \begin{enumerate}[label=(\Alph*)]
        \item $\frac{\partial F}{\partial x}=\frac{\partial F}{\partial y},\frac{\partial^2 F}{\partial x^2}=\frac{\partial^2 F}{\partial y^2}$
        \item $\frac{\partial F}{\partial x}=\frac{\partial F}{\partial y},\frac{\partial^2 F}{\partial x^2}=-\frac{\partial^2 F}{\partial y^2}$
        \item $\frac{\partial F}{\partial x}=-\frac{\partial F}{\partial y},\frac{\partial^2 F}{\partial x^2}=\frac{\partial^2 F}{\partial y^2}$
        \item $\frac{\partial F}{\partial x}=-\frac{\partial F}{\partial y},\frac{\partial^2 F}{\partial x^2}=-\frac{\partial^2 F}{\partial y^2}$
    \end{enumerate}
    \end{multicols}
}

\oneproblem{
    设函数 $u(x,y)=\phi(x+y)+\phi(x-y)+\int_{x-y}^{x+y}\psi(t)dt$，其中函数 $\phi$ 具有二阶导数，$\psi$ 具有一阶导数，则必有（\quad）
    \begin{multicols}{2}
    \begin{enumerate}[label=(\Alph*)]
        \item $\frac{\partial^2 u}{\partial x^2}=-\frac{\partial^2 u}{\partial y^2}$
        \item $\frac{\partial^2 u}{\partial x^2}=\frac{\partial^2 u}{\partial y^2}$
        \item $\frac{\partial^2 u}{\partial x\partial y}=\frac{\partial^2 u}{\partial y^2}$
        \item $\frac{\partial^2 u}{\partial x\partial y}=\frac{\partial^2 u}{\partial x^2}$
    \end{enumerate}
    \end{multicols}
}

\oneproblem{
    设 $f(x,y)=\int_0^{xy}e^{-t^2}dt$，求 $\frac{x}{y}\frac{\partial^2 f}{\partial x^2}-2\frac{\partial^2 f}{\partial x\partial y}+\frac{y}{x}\frac{\partial^2 f}{\partial y^2}$.
}

\oneproblem{
    设函数 $f(u)$ 具有二阶连续导数，而 $z=f(e^x\sin y)$ 满足方程 $\frac{\partial^2 z}{\partial x^2}+\frac{\partial^2 z}{\partial y^2}=e^{2x}z$，求 $f(u)$.
}

\oneproblem{
    已知 $z=u^v,u=\ln\sqrt{x^2+y^2},v=\arctan\frac{y}{x}$，求 $dz$.
}

\oneproblem{
    设函数 $z=f(x,y)$ 在点 $(1,1)$ 处可微，且 $f(1,1)=1$，$f_x'(1,1)=2$，$f_y'(1,1)=3,\varphi(x)=f(x,f(x,x))$，求 $\left.\frac{d}{dx}\varphi^3(x)\right|_{x=1}$.
}

\oneproblem{
    设函数 $f(u)$ 具有二阶连续导数，$z=f(e^x\cos y)$ 满足 $\frac{\partial^2 z}{\partial x^2}+\frac{\partial^2 z}{\partial y^2}=(4z+e^x\cos y)e^{2x}$，若 $f(0)=0,f'(0)=0$，求 $f(u)$ 的表达式.
}

\oneproblem{
    设函数 $f(u)$ 在 $(0,+\infty)$ 内具有二阶导数，且 $z=f(\sqrt{x^2+y^2})$ 满足等式 $\frac{\partial^2 z}{\partial x^2}+\frac{\partial^2 z}{\partial y^2}=0$.
    \begin{enumerate}[label=(\arabic*)]
        \item 验证 $f''(u)+\frac{f'(u)}{u}=0$
        \item 若 $f(1)=0,f'(1)=1$，求函数 $f(u)$ 的表达式.
    \end{enumerate}
}

% ======================================================
% 第 58 天
% ======================================================
\setday{58}{多元复合函数的求导（二）}

\oneproblem{
    设 $w=f(u,v)$ 具有二阶连续偏导数，且 $u=x-cy,v=x+cy$，其中 $c$ 为非零常数，计算并化简 $w_{xx}-\frac{1}{c^2}w_{yy}$.
}

\oneproblem{
    设 $f(x)$ 有连续导数，且 $f(1)=2$。记 $z=f(e^xy^2)$，若 $\frac{\partial z}{\partial x}=z$，求 $f(x)$ 在 $x>0$ 的表达式.
}

\oneproblem{
    设 $f(u,v)$ 具有连续偏导数，且满足 $f_u'(u,v)+f_v'(u,v)=uv$。求 $y(x)=e^{-2x}f(x,x)$ 所满足的一阶微分方程，并求其通解.
}

\oneproblem{
    已知函数 $z=u(x,y)e^{ax+by}$ 且 $\frac{\partial^2 u}{\partial x\partial y}=0$，确定常数 $a,b$，使函数 $z=z(x,y)$ 满足方程 $\frac{\partial^2 z}{\partial x\partial y}-\frac{\partial z}{\partial x}-\frac{\partial z}{\partial y}+z=0$.
}

\oneproblem{
    函数 $u(x,y)$ 具有连续的二阶偏导数，算子 $A$ 定义为 $A(u)=x\frac{\partial u}{\partial x}+y\frac{\partial u}{\partial y}$.
    \begin{enumerate}[label=(\arabic*)]
        \item 求 $A(u-A(u))$；
        \item 利用结论(1)以 $\xi=\frac{y}{x},\eta=x-y$ 为新的自变量改变如下方程形式 $x^2\frac{\partial^2 u}{\partial x^2}+2xy\frac{\partial^2 u}{\partial x\partial y}+y^2\frac{\partial^2 u}{\partial y^2}=0$.
    \end{enumerate}
}

\oneproblem{
    设 $f(t)$ 有二阶连续导数，$r=\sqrt{x^2+y^2}$，$g(x,y)=f\left(\frac{1}{r}\right)$，求 $\frac{\partial^2 g}{\partial x^2}+\frac{\partial^2 g}{\partial y^2}$.
}

\oneproblem{
    设函数 $f(x,y)$ 有二阶连续偏导数，满足 $f_y\neq0$ 且 $f_x^2f_{yy}-2f_xf_yf_{xy}+f_y^2f_{xx}=0$，$y=y(x,z)$ 是由方程 $z=f(x,y)$ 所确定的函数，求 $\frac{\partial^2 y}{\partial x^2}$.
}

\oneproblem{
    已知 $z=xf\left(\frac{y}{x}\right)+2y\varphi\left(\frac{x}{y}\right)$，其中 $f,\varphi$ 均为二阶可微函数.
    \begin{enumerate}[label=(\arabic*)]
        \item 求 $\frac{\partial z}{\partial x},\frac{\partial^2 z}{\partial x\partial y}$；
        \item 当 $f=\varphi$，且 $\left.\frac{\partial^2 z}{\partial x\partial y}\right|_{x=a}=-by^2$ 时，求 $f(y)$.
    \end{enumerate}
}

\oneproblem{
    设二元可微函数 $F(x,y)$ 可写成 $F(x,y)=f(x)+g(y)$，又令 $x=r\cos\theta,y=r\sin\theta$ 时，有 $F(r\cos\theta,r\sin\theta)=S(r)$，试求 $F(x,y)$ 的表达式.
}

\oneproblem{
    设二元函数 $f(x,y)$ 有一阶连续的偏导数，且 $f(0,1)=f(1,0)$。证明：单位圆周上至少存在两点满足方程 $y\frac{\partial}{\partial x}f(x,y)-x\frac{\partial}{\partial y}f(x,y)=0$.
}

\oneproblem{
    设 $F(x_1,x_2,x_3)=\int_0^{2\pi}f(x_1+x_3\cos\varphi,x_2+x_3\sin\varphi)d\varphi$，其中 $f(u,v)$ 具有二阶连续偏导数。已知 $\frac{\partial F}{\partial x_i}=\int_0^{2\pi}\frac{\partial}{\partial x_i}[f(x_1+x_3\cos\varphi,x_2+x_3\sin\varphi)]d\varphi$，$\frac{\partial^2 F}{\partial x_i^2}=\int_0^{2\pi}\frac{\partial^2}{\partial x_i^2}[f(x_1+x_3\cos\varphi,x_2+x_3\sin\varphi)]d\varphi$，$i=1,2,3$，试求 $x_3\left(\frac{\partial^2 F}{\partial x_1^2}+\frac{\partial^2 F}{\partial x_2^2}-\frac{\partial^2 F}{\partial x_3^2}\right)-\frac{\partial F}{\partial x_3}$ 并要求化简.
}

% ======================================================
% 第 59 天
% ======================================================
\setday{59}{隐函数的求导}

\oneproblem{
    设有三元方程 $xy-z\ln y+e^{xz}=1$，根据隐函数存在定理，存在点 $(0,1,1)$ 的一个邻域，在此邻域内该方程（\quad）
    \begin{enumerate}[label=(\Alph*)]
        \item 只能确定一个具有连续偏导数的隐函数 $z=z(x,y)$
        \item 可确定两个具有连续偏导数的隐函数 $x=x(y,z)$ 和 $z=z(x,y)$
        \item 可确定两个具有连续偏导数的隐函数 $y=y(x,z)$ 和 $z=z(x,y)$
        \item 可确定两个具有连续偏导数的隐函数 $x=x(y,z)$ 和 $y=y(x,z)$
    \end{enumerate}
}

\oneproblem{
    设函数 $z=z(x,y)$ 由方程 $F(x-y,z)=0$ 确定，其中 $F(u,v)$ 具有连续二阶偏导数，则 $\frac{\partial^2 z}{\partial x\partial y}$.
}

\oneproblem{
    设 $z=z(x,y)$ 是由方程 $2\sin(x+2y-3z)=x+2y-3z$ 所确定的二元隐函数，求 $\frac{\partial z}{\partial x}+\frac{\partial z}{\partial y}$.
}

\oneproblem{
    设 $z=z(x,y)$ 由方程 $F\left(x+\frac{z}{y},y+\frac{z}{x}\right)=0$ 所决定，其中 $F(u,v)$ 具有连续偏导数，且 $xF_u+yF_v\neq0$，求 $x\frac{\partial z}{\partial x}+y\frac{\partial z}{\partial y}$，结果要求不显含有 $F$ 及其偏导数.
}

\oneproblem{
    设函数 $f(x,y)$ 有二阶连续偏导数，满足 $f_y\neq0$ 且 $f_x^2f_{yy}-2f_xf_yf_{xy}+f_y^2f_{xx}=0$，$y=y(x,z)$ 是由方程 $z=f(x,y)$ 所确定的函数，求 $\frac{\partial^2 y}{\partial x^2}$.
}

\oneproblem{
    设 $z=z(x,y)$ 是由方程 $F\left(z+\frac{1}{x},z-\frac{1}{y}\right)=0$ 确定的隐函数，且具有连续的二阶偏导数，求证：$x^2\frac{\partial z}{\partial x}-y^2\frac{\partial z}{\partial y}=1$ 和 $x^3\frac{\partial^2 z}{\partial x^2}+xy(x-y)\frac{\partial^2 z}{\partial x\partial y}-y^3\frac{\partial^2 z}{\partial y^2}+2=0$.
}

\oneproblem{
    设函数 $f(u,v)$ 可微，$z=z(x,y)$ 由方程 $(x+1)z-y^2=x^2f(x-z,y)$ 确定，求 $dz|_{(0,1)}$.
}

\oneproblem{
    设 $f(x,y,z)=e^xyz^2$，其中 $z=z(x,y)$ 是由 $x+y+z+xyz=0$ 确定的隐函数，求 $f_x'(0,1,-1)$.
}

\oneproblem{
    已知 $xy=xf(z)+yg(z),xf'(z)+yg'(z)\neq0$，其中 $z=z(x,y)$ 是 $x$ 和 $y$ 的函数. 求证：$[x-g(z)]\frac{\partial z}{\partial x}=[y-f(z)]\frac{\partial z}{\partial y}$.
}

\oneproblem{
    设 $u=f(x,y,z),\varphi(x^2,e^y,z)=0,y=\sin x$，其中 $f,\varphi$ 都具有一阶连续偏导数，且 $\frac{\partial\varphi}{\partial z}\neq0$，求 $\frac{du}{dx}$.
}

\oneproblem{
    设 $u=f(x,y,z)$ 有连续偏导数，$y=y(x)$ 和 $z=z(x)$ 分别由方程 $e^{xy}-y=0$ 和 $e^z-xz=0$ 所确定，求 $\frac{du}{dx}$.
}

\oneproblem{
    设 $u=f(x,y,z)$ 有连续的一阶偏导数，又函数 $y=y(x)$ 及 $z=z(x)$ 分别由下列两式确定：$e^{xy}-xy=2$ 和 $e^x=\int_0^{x-z}\frac{\sin t}{t}dt$，求 $\frac{du}{dx}$.
}

\oneproblem{
    设 $z=z(x,y)$ 是由方程 $x^2+y^2-z=\varphi(x+y+z)$ 所确定的函数，其中 $\varphi$ 具有2阶导数且 $\varphi'\neq-1$.
    \begin{enumerate}[label=(\arabic*)]
        \item 求 $dz$
        \item 记 $u(x,y)=\frac{1}{x-y}\left(\frac{\partial z}{\partial x}-\frac{\partial z}{\partial y}\right)$，求 $\frac{\partial u}{\partial x}$.
    \end{enumerate}
}

% ======================================================
% 第 60 天
% ======================================================
\setday{60}{多元函数微分学的几何应用}

\oneproblem{
    设 $f(x,y)$ 在点 $(0,0)$ 附近有定义，且 $f_x'(0,0)=3$，$f_y'(0,0)=1$，则（\quad）
    \begin{enumerate}[label=(\Alph*)]
        \item $dz|_{(0,0)}=3dx+dy$
        \item 曲面 $z=f(x,y)$ 在 $(0,0,f(0,0))$ 处的法向量为 $(3,1,1)$
        \item 曲线 $\begin{cases}z=f(x,y)\\y=0\end{cases}$ 在 $(0,0,f(0,0))$ 处的切向量为 $(1,0,3)$
        \item 曲线 $\begin{cases}z=f(x,y)\\y=0\end{cases}$ 在 $(0,0,f(0,0))$ 处的切向量为 $(3,0,1)$
    \end{enumerate}
}

\oneproblem{
    过点 $(1,0,0)$ 与 $(0,1,0)$ 且与 $z=x^2+y^2$ 相切的平面方程为（\quad）.
    \begin{multicols}{2}
    \begin{enumerate}[label=(\Alph*)]
        \item $z=0$ 与 $x+y-z=1$
        \item $z=0$ 与 $2x+2y-z=2$
        \item $y=x$ 与 $x+y-z=1$
        \item $y=x$ 与 $2x+2y-z=2$
    \end{enumerate}
    \end{multicols}
}

\oneproblem{
    求曲面 $z=\frac{x^2}{2}+y^2$ 平行于平面 $2x+2y-z=0$ 的切平面方程.
}

\oneproblem{
    试求由曲线 $\begin{cases}3x^2+2y^2=12\\z=0\end{cases}$ 绕 $y$ 轴旋转一周得到的旋转面在点 $(0,\sqrt{3},\sqrt{2})$ 处的指向外侧的单位法向量.
}

\oneproblem{
    设 $a,b,c,\mu>0$，曲面 $xyz=\mu$ 与曲面 $\frac{x^2}{a^2}+\frac{y^2}{b^2}+\frac{z^2}{c^2}=1$ 相切，求 $\mu$ 的值.
}

\oneproblem{
    求椭球面 $x^2+2y^2+3z^2=21$ 上某点 $M$ 处的切平面 $\pi$ 的方程，使平面 $\pi$ 过已知直线 $L:\frac{x-6}{2}=\frac{y-3}{1}=\frac{2z-1}{-2}$.
}

\oneproblem{
    设 $L:\begin{cases}x+y+b=0\\x+ay-z-3=0\end{cases}$ 在平面 $\pi$ 上，而平面 $\pi$ 与曲面 $z=x^2+y^2$ 相切于点 $(1,-2,5)$，求 $a,b$ 之值.
}

\oneproblem{
    过直线 $\begin{cases}10x+2y-2z=27,\\x+y-z=0\end{cases}$ 作曲面 $3x^2+y^2-z^2=27$ 的切平面，求此切平面的方程.
}

\oneproblem{
    证明曲面 $F(x-my,z-ny)=0$ 的所有切平面恒与定直线平行，其中 $F(u,v)$ 可微.
}

\oneproblem{
    证明：曲面 $f\left(\frac{x-a}{z-c},\frac{y-b}{z-c}\right)=0$ 上任一点处的切平面过某一定点，其中 $f(u,v)$ 为可微函数.
}

\oneproblem{
    证明曲面 $z+\sqrt{x^2+y^2+z^2}=x^3f\left(\frac{y}{x}\right)$ 上任意点处的切平面在 $Oz$ 轴上的截距与切点到坐标原点的距离之比为常数，并求此常数.
}

\oneproblem{
    设 $S$ 为 $\mathbb{R}^3$ 中的抛物面 $z=\frac{1}{2}(x^2+y^2)$，$P=(a,b,c)$ 为 $S$ 外一固定点，满足 $a^2+b^2>2c$。过 $P$ 作 $S$ 的所有切线。证明：这些切线的切点落在同一张平面上.
}

\oneproblem{
    设 $f(u,v)$ 在全平面上有连续的偏导数，证明：曲面 $f\left(\frac{x-a}{z-c},\frac{y-b}{z-c}\right)=0$ 的所有切平面都交于点 $(a,b,c)$.
}

\oneproblem{
    设 $F(x,y,z),G(x,y,z)$ 有连续偏导数，$\frac{\partial(F,G)}{\partial(x,z)}\neq0$，曲线 $\begin{cases}F(x,y,z)=0,\\G(x,y,z)=0\end{cases}$ 过点 $P_0(x_0,y_0,z_0)$。记 $\Gamma$ 在 $xOy$ 面上的投影曲线为 $S$。求 $S$ 上过点 $(x_0,y_0)$ 的切线方程.
}

\oneproblem{
    已知曲面 $e^{2x-z}=f(\pi y-\sqrt{2}z)$，且 $f$ 可微，证明该曲面为柱面.
}